% ===================================================================
%  اللوحة نمرة ١٠ — البحر. الميناء.
%
%  Traduction, faite en 2026, en arabe egyptien ecrit, sur l'ido de
%  texto/io. Regles de la colonne : voir 10-lawha-01.tex. Une seule
%  numerotation.
%
%  C'EST LE TABLEAU OU L'EGYPTE EST LA CHOSE DECRITE. Le pilote, la
%  drague, l'ecluse, le canal, l'isthme, le phare, l'arsenal, le
%  tramway du port : ce ne sont pas des mots qu'il a fallu chercher,
%  ce sont ceux qu'on lit sur les avis de navigation d'Alexandrie,
%  de Port-Said et de Suez.
%
%    مرشد (36)     le pilote      — le titre officiel du pilote du canal
%    كرّاكة (68)    la drague      — celles qui ont creuse le canal
%    هويس (69)     l'ecluse
%    قنال (67)     le canal       — en Egypte, LE canal
%    برزخ (90)     l'isthme       — برزخ السويس
%    ترسانة (18)   le chantier naval
%    فنار (93)     le phare       — le premier etait a Alexandrie
%    ترام (74)     le tramway     — celui d'Alexandrie
%    جمرك (65)     la douane
%    ونش (32, 62)  la grue
%    صندل (17, 41) la peniche, le chaland
%    رفّاص (28)     l'helice
%    لنش (16, 71)  la vedette, le remorqueur
%
%  ET LE MOT « وابور » COUVRE TOUT CE QUI MARCHE A LA VAPEUR. Il a
%  deja servi au tableau 8 pour la locomobile de la batteuse ; il
%  sert ici pour le paquebot (3). En Egypte il fait aussi la
%  locomotive et le fourneau. Un seul mot pour une seule invention,
%  la ou le francais en a quatre.
%
%  LES GRADES DE MARINE SONT TURCS, COMME LES METIERS EN -جي :
%
%    يوزباشي (79, 80)  le capitaine
%    باشجاويش (85)     le quartier-maitre
%    قومندان (83)      le commandant
%
%  Le turc a laisse a l'egyptien son armee et ses ateliers ; l'italien
%  lui a laisse sa maison et ses vetements. Les deux couches se lisent
%  encore, tableau par tableau.
% ===================================================================

%%K t10-apar-1 apar
\VUsaut{4.0mm}
\VUtitre{15.0pt}{60}{اللوحة نمرة ١٠}
\VUsaut{3.0mm}
\VUpk{11.4pt}{40}{البحر. --- الميناء.}
\VUsaut{3.0mm}

%%K t10-01-1 p
1. --- في الصيف، في ناس بتدوّر على الهدوء والبرد في الجبل، وناس تانية بتفضّل
تروح البحر. وكده عملنا السنة اللي فاتت، علشان إحنا بنحبّ \VUgras{المحيط}
\textsuperscript{(1)} أوي.

%%K t10-02-1 p
2. --- وأنا شخصيًا بحسّ بمتعة كبيرة وأنا بتأمّل المساحة المية الهايلة دي اللي
بتمتدّ لحدّ \VUgras{الأفق} \textsuperscript{(2)}، وأنا بشوف \VUgras{الوابورات}
\textsuperscript{(3)} الضخمة بتسافر \VUgras{سفريات} طويلة، شايلة المسافرين الرايحين
لأمريكا.

%%K t10-03-1 p
3. --- علشان كده عمّي، اللي عارف ميلي ده، بيختار دايمًا لقضاء الأجازة شاطئ هادي
أوي، جنب واحد من \VUgras{موانينا الحربية} الكبيرة.

%%K t10-03-2 p
وفي آخر إقامة لنا، \VUgras{الأسطول} جه رسا ورا \VUgras{حاجز الأمواج}
\textsuperscript{(4)} الضخم اللي بيقفل ويحمي \VUgras{الميناء الحربي}
\textsuperscript{(5)}، وقدرت أتفرّج زي ما كنت عايز على \VUgras{السفن الحربية}
المختلفة، و\VUgras{الغوّاصة} \textsuperscript{(10)} الصغيرة اللي بتغطس وتختفي تحت
المية، وكمان \VUgras{المدرّعة} \textsuperscript{(9)} الضخمة اللي، لحدّ دلوقت، ما
كانتش بتخاف تقريبًا غير من هجمات \VUgras{طرّاد الطوربيد} \textsuperscript{(8)}.

%%K t10-tit-2 sub
\VUsaut{3.0mm}
\VUpk{10.2pt}{40}{المراكب الصغيرة.}
\VUsaut{3.0mm}

%%K t10-04-1 p
4. --- وده كان \VUgras{بيعرف} من زمان يلاقي النقطة الضعيفة في الدرع المعدني
التخين علشان يحطّ فيها \VUgras{الطوربيد} الخطر، اللي انفجاره بيغرّق المركب؛ بس
مع كده كان أقل خوفًا من الغوّاصة، علشان مدمّرات الطوربيد بتلحقه بسهولة.

%%K t10-05-1 p
5. --- وقدرت كمان أشوف \VUgras{الطرّادات} \textsuperscript{(6)} السريعة المدرّعة،
اللي تقريبًا صعب اختراقها زي المدرّعة الحقيقية؛ وكام \VUgras{سفينة نقل}
\textsuperscript{(12)}؛ وتلاتة أو أربع \VUgras{زوارق مدافع} \textsuperscript{(7)}؛
وأخيرًا \VUgras{فرقاطة} \textsuperscript{(11)} واحدة.
\VUgras{ومنظر الأسطول ده وهو بيناور علشان يقلع أحلى منظر ممكن.}

%%K t10-06-1 p
6. --- وقدّ إيه الفرق في البناء بين أكوام الصلب الكبيرة دي و\VUgras{المراكب أبو
تلات صواري} الأنيقة أو \VUgras{المراكب الشراعية} \textsuperscript{(13)} الرشيقة اللي
لسه بتتستعمل في الأسطول التجاري، اللي شفتها داخلة الميناء بكل شراعها وأنا
بتمشّى على \VUgras{الرصيف} \textsuperscript{(14)}. والحقيقة إن دي بتفقد كتير من
مميزاتها \VUgras{لما} \VUgras{الريح} يبقى معاكس. لازم ينزّلوا كل الشراع علشان
يدخلوا الحوض تاني، ولازم \VUgras{لنش قاطر} \textsuperscript{(16)} علشان يجيبهم
ويودّيهم، زي \VUgras{الصنادل} \textsuperscript{(17)} العادية، للرصيف.

%%K t10-tit-3 sub
\VUsaut{3.0mm}
\VUpk{10.2pt}{40}{الميناء التجاري.}
\VUsaut{3.0mm}

%%K t10-07-1 p
7. --- واللي مشوّق في الميناء اللي بتكلّم عنه، إنه مش بس فيه ميناء حربي معمول
بالإيد بأشغال عملاقة، لكن فيه كمان \VUgras{ميناء تجاري} رائع في \VUgras{مصبّ}
نهر كبير.

%%K t10-08-1 p
8. --- ولسان الأرض اللي بيفصل الحوضين متحصّن ومحمي بسلسلة \VUgras{بطاريات}
و\VUgras{حصون} بتتقدّم في البحر. ولازم تصريح مخصوص علشان تتمشّى على
\VUgras{الأسوار} \textsuperscript{(15)}. وأنا خدته بسهولة عن طريق خالي، اللي مهندس في
\VUgras{الترسانة} \textsuperscript{(18)}،

%%K t10-08-1 p suite
ورحت أزور المراكب اللي بيبنوها تحت هناجر واسعة.

%%K t10-09-1 p
9. --- وحتى حضرت إنزال مدرّعة، وفاكر اللحظة اللي حسّ بيها الجمهور كله بالانفعال
لما الكتلة الضخمة، لما سابوها لنفسها، ابتدت تزحلق على هيكلها المهدهد وجت تعوم
على مية النهر.

%%K t10-10-1 p
10. --- وعلشان توصل للترسانة، بتعدّي \VUgras{ميدان التدريب} \textsuperscript{(19)}
اللي عساكر الحامية بيتمرّنوا فيه كل يوم، وبتعدّي على \VUgras{كوبري حديد}
\textsuperscript{(20)} صغير، اللي بيوصّل ميدان العرض بـ\VUgras{الأرسنال}
\textsuperscript{(21)}.

%%K t10-tit-4 sub
\VUsaut{3.0mm}
\VUpk{10.2pt}{40}{الأرسنال والأحواض.}
\VUsaut{3.0mm}

%%K t10-11-1 p
11. --- ورحت مع خالي أزور المبنى الكبير ده المليان \VUgras{مدافع}
\textsuperscript{(22)}، و\VUgras{قنابل} \textsuperscript{(23)}، و\VUgras{شظايا}. وشفت
كمان \VUgras{مصنع المعادن} \textsuperscript{(24)} اللي بتتعمل فيه \VUgras{غلايات}
\textsuperscript{(25)} الوابورات.

%%K t10-12-1 p
12. --- وقريّب أوي، في \VUgras{الأحواض} \textsuperscript{(26)} --- أحواض التصليح ---
و\VUgras{الأحواض العايمة} \textsuperscript{(27)} اللي تستاهل النظر أوي، واللي قدرت
أشوف فيها من قريب أوي، على مركب تحت التصليح، \VUgras{الرفّاص} \textsuperscript{(28)}،
و\VUgras{القاعدة} \textsuperscript{(29)}، و\VUgras{الدفّة} \textsuperscript{(30)} بتاعت
المركب.

%%K t10-13-1 p
13. --- والميناء التجاري يستاهل الدراسة زي الميناء الحربي بالظبط، وقضّيت ساعات
كتير بتفرّج على الأرصفة اللي راسية عليها \VUgras{البواخر أبو عجلة}
\textsuperscript{(31)} اللي بتطلع عكس النهر، وباتفرّج على \VUgras{الأوناش}
\textsuperscript{(32)} وهي بتشتغل، اللي بترفع، زي الريشة، حمولات ضخمة.

%%K t10-tit-5 sub
\VUsaut{4.0mm}
\VUpk{10.2pt}{40}{المرشد.}
\VUsaut{3.0mm}

%%K t10-14-1 p
14. --- واتفرّجت على \VUgras{بواخر الأطلنطي} \textsuperscript{(34)} وهي خارجة من
المرسى، و\VUgras{أعلامها} \textsuperscript{(35)} مفرودة في الهوا، وبيقودها
\VUgras{مرشد} \textsuperscript{(36)} شاطر، اللي بيعرف بمهارة يمشي في \VUgras{الممر
الضيّق} المعلّم بـ\VUgras{العوّامات} \textsuperscript{(40)} و\VUgras{الشمندورات}،
وهو بيتفادى \VUgras{الصنادل} \textsuperscript{(41)} اللي بتتساق بصعوبة بشراعها
المربّع الواحد. ومّيزت على \VUgras{المقدّمة} \textsuperscript{(37)}
\VUgras{الكبائن} \textsuperscript{(39)} بتاعة الدرجة التانية، وعلى
\VUgras{المؤخّرة} \textsuperscript{(38)} صالات ركّاب الدرجة الأولى.

%%K t10-15-1 p
15. --- وأحيانًا حضرت كمان شحن \VUgras{صندل} \textsuperscript{(42)} اتكوّمت فيه حالًا
بضايع \VUgras{المخزن} \textsuperscript{(43)} و\VUgras{محطة الميناء} \textsuperscript{(44)}،
اللي جابتها قطارات كتير أوي بتاعة \VUgras{سكة الرصيف} \textsuperscript{(45)}.

%%K t10-tit-6 sub
\VUsaut{3.0mm}
\VUpk{10.2pt}{40}{متعة التجديف.}
\VUsaut{3.0mm}

%%K t10-16-1 p
16. --- وكتير جدّفت في \VUgras{الحوض العايم} \textsuperscript{(53)}، اللي بتلجأ له
\VUgras{القوارب} \textsuperscript{(46)}، وكانت متعة لي إني أمسك \VUgras{المجداف}
\textsuperscript{(47)}. وطلعت \VUgras{سطح} \textsuperscript{(48)} \VUgras{مركب شراعي}
\textsuperscript{(49)}، وتسلّقت، بـ\VUgras{الحبال} \textsuperscript{(52)}، على
\VUgras{الصاري} \textsuperscript{(51)} لحدّ \VUgras{العوارض} \textsuperscript{(50)}؛
وبعدين كمّلت نزهتي \VUgras{بقارب صغير} \textsuperscript{(54)} بين \VUgras{مراكب
الصيد} \textsuperscript{(55)} التقيلة، اللي \VUgras{الشباك} \textsuperscript{(56)}
بتتنشّف فيها.

%%K t10-17-1 p
17. --- وعلى \VUgras{الرصيف} \textsuperscript{(58)} عدّت قدّامي \VUgras{بضايع}
\textsuperscript{(59)} من كل نوع، محطوطة في \VUgras{صناديق} \textsuperscript{(60)} أو
في \VUgras{بالات} \textsuperscript{(61)}. ودي كانت \VUgras{حمولة} \textsuperscript{(63)}
الصنادل، و\VUgras{أوناش} \textsuperscript{(62)} قوية بترفعها وتنزّلها على
\VUgras{لواري} \textsuperscript{(64)}، و\VUgras{أصحاب النقل} بيقدّموا الإقرار
ويدفعوا الرسوم في \VUgras{الجمرك} \textsuperscript{(65)}.

%%K t10-18-1 p
18. --- وأخيرًا، لما وصلت لـ\VUgras{الكوبري الدوّار} \textsuperscript{(66)} أو
لـ\VUgras{الهويس} \textsuperscript{(69)}، وأنا بعدّي قدام \VUgras{الكرّاكة}
\textsuperscript{(68)} اللي بتنضّف \VUgras{القنال} \textsuperscript{(67)}، نزلت على
الرصيف جنب \VUgras{برج الإشارات} \textsuperscript{(70)}.

%%K t10-19-1 p
19. --- وعلى الناحية التانية من الرصيف ده بترسي \VUgras{اللنشات}
\textsuperscript{(71)} اللي بتنزّل ضبّاط الأسطول بر. ومنظر البحّارة وهما بيرفعوا
مجاديفهم على وزن واحد يستاهل الفرجة، والقارب اللي شايلاه \VUgras{الموجة}
\textsuperscript{(72)} بيوصل بسهولة لسلّم النزول. ومن المكان ده، الواحد يقدر يلاحظ
أحسن \VUgras{المدّ والجزر} وحركة \VUgras{الجزر} و\VUgras{المدّ} على
\VUgras{الشاطئ} \textsuperscript{(73)}.

%%K t10-tit-7 sub
\VUsaut{3.0mm}
\VUpk{10.2pt}{40}{النادي.}
\VUsaut{3.0mm}

%%K t10-20-1 p
20. --- وعلشان أرجع البيت، ركبت \VUgras{الترام} \textsuperscript{(74)} اللي
\VUgras{محطته} \textsuperscript{(75)} جنب \VUgras{العمود} اللي قدام نادي الضبّاط.

%%K t10-20-2 p
ومن \VUgras{بلكونة} \textsuperscript{(76)} النادي، المزيّنة بـ\VUgras{مراسي}
\textsuperscript{(77)} ومدافع وقنابل، كتير شفت تجمّع رائع لضبّاط البحرية:
\VUgras{ملازمين} \textsuperscript{(78)} بسيوفهم، و\VUgras{يوزباشية مدفعية}
\textsuperscript{(79)} و\VUgras{مشاة} \textsuperscript{(80)} بـ\VUgras{سيوفهم}.
وورا منهم كان في \VUgras{بحّار} \textsuperscript{(81)} بيناولهم \VUgras{النضّارة
المكبّرة} لما يبقوا عايزين يشوفوا اللي بيحصل بعيد.

%%K t10-21-1 p
21. --- وشفت كمان \VUgras{ملازمين تانيين} \textsuperscript{(82)} صغيّرين بياخدوا
تعليمات \VUgras{القومندان} \textsuperscript{(83)} بتاعهم أو \VUgras{الأميرال}
\textsuperscript{(84)}، والـ\VUgras{باشجاويش} \textsuperscript{(85)} مستنّي علشان
يوديها للمركب.

%%K t10-22-1 p
22. --- ومن البلكونة دي، المنظر رائع: الواحد بيشرف على الشاطئ كله
بـ\VUgras{خيامه}

%%K t10-22-1 p suite
\textsuperscript{(86)} و\VUgras{كبائنه} \textsuperscript{(87)}، وبيشوف، في ساعة
الاستحمام، \VUgras{المستحمّين} \textsuperscript{(88)} بيلعبوا بفرحة قدام
\VUgras{الكازينو} \textsuperscript{(89)}.

%%K t10-23-1 p
23. --- وفي آخر الشاطئ، الساحل بيعمل \VUgras{شبه جزيرة} \textsuperscript{(91)} صغيرة
صخرية، اللي \VUgras{الموج الكبير} \textsuperscript{(92)} بيجي يتكسّر عليها، واللي
اتبنى عليها \VUgras{فنار} \textsuperscript{(93)}، بيدلّ الملاحين على الطريق بالليل.
وشبه الجزيرة دي متوصّلة بالبرّ بـ\VUgras{برزخ} ضيّق أوي \textsuperscript{(90)} بس.

%%K t10-tit-8 sub
\VUsaut{3.0mm}
\VUpk{10.2pt}{40}{منظر عام للسواحل.}
\VUsaut{3.0mm}

%%K t10-24-1 p
24. --- و\VUgras{الجرف} \textsuperscript{(94)} بيمتدّ، مشقوق من البحر، اللي بيرسم فيه على
مزاجه سلسلة \VUgras{خلجان} \textsuperscript{(95)} و\VUgras{رؤوس}، وبيخلّيه حلو زي
الصورة.

%%K t10-24-2 p
وعلى \VUgras{الهضبة} \textsuperscript{(96)}، اللي بتشرف على \VUgras{المضيق}
\textsuperscript{(98)}، في \VUgras{طابية} \textsuperscript{(97)} مهمة أوي وكام «فيلا».

%%K t10-25-1 p
25. --- والمضيق ده بيفصل عن القارّة \VUgras{جزيرة} \textsuperscript{(99)} خطرة أوي،
محوّطة بـ\VUgras{شعاب} \textsuperscript{(100)} و\VUgras{أمواج متكسّرة}، اللي بتجنح
عليها أحيانًا قوارب وحتى مراكب كبيرة. وبالرغم من سرعة النجدة وسرعة وصول
\VUgras{قوارب النجاة}، إلا إن \VUgras{الغرق} ده بيرعب، وفي \VUgras{عواصف}
الاعتدال غرقت كام مركب بناسها وبضايعها.

%%K t10-25-2 p
وبالرغم من الجيرة دي، \VUgras{الميناء} البحّارة بيترددوا عليه أوي.

\VUsaut{6.0mm}
\VUfilet{22mm}
